\chapter{Proposal Feedback and Iteration Plan}

\section{Proposal Feedback Summary}
The proposal defense result was issued by the FYP Committee by email on 14~September~2026 (Figure~\ref{fig:panelemail}). The decision was \textbf{Approved with Minor Modifications}. The committee judged the idea sound and told the team to focus on implementation, but asked for three things in particular: clarity on how the system would handle scale, drift, incomplete visual coverage, lighting variation, contrast, low-texture surfaces and reconstruction quality; a clearly defined practical capture scope and set of operating conditions, including the type and size of sites targeted and the limitations of commodity Android devices; and capture guidance and the handling of poor-quality input treated as part of the solution rather than left to the user. The committee also noted that the scope looked large, that feasibility should be checked deeply, and that the presentation should have stated plainly that the scope is limited to budget phone video.
\section{Proposal Action Register}
\label{sec:actionregister}
Table~\ref{tab:actionregister} records each point the committee raised and what was done about it.

{\small
\begin{longtable}{p{0.03\textwidth}p{0.20\textwidth}p{0.42\textwidth}p{0.07\textwidth}p{0.16\textwidth}}
\caption{Proposal action register. Feedback is condensed from the committee's email of 14~September~2026 (Figure~\ref{fig:panelemail}).}\label{tab:actionregister}\\
\toprule
\textbf{No.} & \textbf{Committee feedback} & \textbf{Action taken} & \textbf{Status} & \textbf{Evidence} \\
\midrule
\endfirsthead
\toprule
\textbf{No.} & \textbf{Committee feedback} & \textbf{Action taken} & \textbf{Status} & \textbf{Evidence} \\
\midrule
\endhead
1 & Handling of scale & Sim3 alignment of the COLMAP reconstruction to ARCore's metric frame, validated against a tape measure: mean error 4.8\,mm (0.19\%), worst case 0.29\%. & Done & Figure~\ref{fig:sim3}, Table~\ref{tab:tape} \\
\midrule
2 & Handling of drift & Diagnosed ARCore's silent re-anchoring (about 1\,m and 28\textdegree) as the cause of ghosting, and replaced ARCore poses with image-based COLMAP poses. Anchor logging makes re-anchor events detectable, and the IMU bridges tracking gaps: 3.8\,cm error over a simulated 9\,s gap, against 30.6\,cm for interpolation. & Done & Table~\ref{tab:imufallback} \\
\midrule
3 & Incomplete visual coverage & Live coverage overlay: depth is accumulated into 10\,cm voxels, each recording which of 16 directions it was seen from, drawn back as red, yellow and green dots at 2\,Hz while recording. & Done & Figure~\ref{fig:coverage}, Figure~\ref{fig:appcapture} \\
\midrule
4 & Lighting variation and contrast & Live IQA warns on darkness during capture, and training applies per-image bilateral-grid exposure compensation for the phone's auto-exposure. Capture across markedly different lighting conditions has not yet been tested. & Partial & Section~\ref{sec:completed}, Table~\ref{tab:risks} \\
\midrule
5 & Low-texture surfaces & Live IQA warns on low texture, and the IMU fallback bridges the tracking loss that plain walls cause. Reconstruction on large untextured surfaces remains weaker than on textured ones. & Partial & Table~\ref{tab:imufallback}, Chapter~\ref{ch:testing} \\
\midrule
6 & Reconstruction quality & Acceptance thresholds fixed at PSNR $\geq$ 25\,dB and SSIM $\geq$ 0.85 on held-out views. The current best model (v6) reaches 26.52\,dB and 0.8886 on 42 held-out frames. & Done & Table~\ref{tab:quality} \\
\midrule
7 & Define the practical capture scope and operating conditions, including site size and the limits of commodity Android & Scope now states single rooms or a connected set of rooms, captured on a commodity Android phone (realme, no LiDAR), with the lighting and tracking conditions ARCore needs listed as explicit constraints. & Done & Sections~1.6 and~1.10 \\
\midrule
8 & Capture guidance and handling of poor-quality input should be part of the solution & Live IQA for blur, brightness, motion speed and texture, with a colour-coded banner and a vibration alert, plus the coverage overlay; poor frames are then dropped at preprocessing by gyro-aware sharp-frame selection. & Done & Section~\ref{sec:completed}, Figure~\ref{fig:appcapture} \\
\midrule
9 & State plainly that the scope is limited to budget phone video & The abstract and the scope section both state that capture is from a commodity Android phone with no LiDAR and no special hardware. & Done & Abstract, Section~1.6 \\
\midrule
10 & Work on edge devices and handle every edge case & Edge and failure cases are documented with their mitigations: tracking loss, plain walls, mirrors and glass, network interruption and pod failure. Mirrors, glass and large rooms remain open and are carried in the risk register. & Partial & Chapter~\ref{ch:testing}, Table~\ref{tab:risks} \\
\midrule
11 & Check feasibility deeply; focus on implementation & The end-to-end pipeline is implemented and measured on real captures, from phone recording through cloud GPU reconstruction to a metric model. Remaining work is listed honestly rather than claimed. & Done & Chapter~\ref{ch:evidence}, Chapter~\ref{ch:testing} \\
\bottomrule
\end{longtable}
}
\section{Iteration Strategy}
We undertake a strategy that comprises two iterations, represented by FYP-1 and FYP-2:

\begin{itemize}
    \item \textbf{Iteration~1} (FYP-1, Weeks 1--8): Implementation of the key model; thus, we create the app, reconstruction pipeline and processing server, integrate cloud GPU, and test everything on the real data. As of the mid-evaluation, the core pipeline is complete and refinement is ongoing. The deliverables include a functioning android app, FastAPI server, Run Pod integration, and accurate 3D reconstruction.
    \item \textbf{Iteration~2} (FYP-2, Weeks 9--16): Performance tuning and enrichment of the functions offered. The deliverables include the reduction of time necessary for processing and improvement of robustness of the pipeline; creation of the web-viewing page based on Three.js, which can also measure; the structural layer with floor plans; and field validation on multiple rooms.
\end{itemize}

\section{Iteration 1 Plan}
\begin{itemize}
    \item \textbf{Objective}: Create the complete system from mobile phone capture to 3D model.
    \item \textbf{Features}: ARCore capture, including IQA \& overlay capabilities; COLMAP SfM together with splatfacto training; FastAPI server with RunPod dispatch; upload interface inside the Android application.
    \item \textbf{Deliverables}: Working Android APK, processing server, cloud GPU integration, verification of reconstruction quality in two rooms.
    \item \textbf{Evidence}: Live demo of mobile phone capture $\rightarrow$ upload $\rightarrow$ cloud processing $\rightarrow$ 3D model display.
\end{itemize}

\section{Iteration 2 Plan}
\begin{itemize}
    \item \textbf{Aim}: Speed up processing time, enhance pipeline stability, finalize the structural layer, and integrate measuring instruments.
    \item \textbf{Characteristics}: pipeline performance enhancement (goal $<$35 minutes on RTX 4090, projected; 75 minutes measured on RTX 3070), edge case resilient, wall-snapping measurement using the web viewer, and completion of the structural layer; modules and timetable in Chapter~\ref{ch:plan}.
    \item \textbf{Expected results}: Optimized pipeline with lower processing time; web viewer providing measurement tools; structural layer capable of floor plan export; report on field validation for $\geq$5 rooms.
    \item \textbf{Proof}: Processing time benchmarks; measuring accuracy comparison with tape measure on at least 5 rooms.
\end{itemize}

\section{Work Breakdown and Responsibilities}

Table~\ref{tab:workbreakdown} records who owns each task and its status at the mid-evaluation.

\begin{table}[htbp]
\centering
\small
\caption{Work breakdown and ownership}
\label{tab:workbreakdown}
\begin{tabular}{p{4cm}p{3.5cm}p{3cm}}
\toprule
\textbf{Task} & \textbf{Owner} & \textbf{Status} \\
\midrule
ARCore capture app & Tayyab, Umair & Complete \\
Live IQA + coverage overlay & Tayyab (IQA), Umair and Tayyab (coverage) & Complete \\
Upload UI (ProcessingActivity) & Umair & Complete \\
Reconstruction pipeline (SfM) & Tayyab & Core complete; refinement ongoing \\
3DGS training \& analysis & Tayyab; trainer fixes by Zaid & Refinement ongoing \\
FastAPI processing server & Umair & Complete \\
RunPod cloud GPU integration & Umair & Complete \\
CI/CD (GitHub Actions) & Zaid, Umair & Complete \\
Unit tests (pytest) & Zaid & Complete \\
Pipeline optimization (FYP-2) & Umair, Tayyab & Planned \\
Web viewer (FYP-2) & Zaid, Tayyab & Planned \\
3D mesh creation & Team & Planned; not yet scheduled \\
Structural layer (\texttt{structure.py}) & Tayyab & Prototype; completion FYP-2 \\
\bottomrule
\end{tabular}
\end{table}

\section{Progress Against Plan}
The core-pipeline milestones for Iteration~1 have been accomplished; refining work is still going on:
\begin{itemize}
    \item Weeks 1--2: ARCore capture app with primitive recording feature has been achieved.
    \item Weeks 3--4: COLMAP SfM pipeline with ARCore Sim3 alignment has been completed. The main issues leading to low quality of reconstruction (gaps in tracking, distortion of lens, and problems with splats) have been solved.
    \item Weeks 5--6: Live IQA has been realized with coverage overlay and a number of improvements in quality (selection of gyro-aware frames, MVS initialization).
    \item Weeks 7--8: the processing server has been developed and integrated with RunPod, the upload interface has been built, and CI/CD has been finalized.
    \item \textbf{Current issues to be addressed}: 3DGS refinement, 3D mesh generation, finishing the structural layer and the web workspace.
\end{itemize}

\section{Delays, Deviations and Recovery Plan}
\begin{itemize}
    \item \textbf{Deviation}: At first, we came up with our unique script for training (\texttt{train.py}) that included depth supervision, flattening regularization, and pose correction. But after carrying out a thorough analysis (detailed in \texttt{Training\_Analysis.md}) we realized that it had proved less effective than splatfacto. Consequently, we switched to splatfacto, and kept our unique preprocessing intact. The overall process took around two weeks, but the analysis allowed us to make the right decision about the design.
    \item \textbf{Recovery}: The basic pipeline has been launched, but Iteration~1 is not complete yet: there are still some points left to resolve. Therefore, the FYP-2 project will make use of the one-command v6 pipeline and the processing time decrease, and will make plans for structural layer and web workspace in Iteration~2a-2b. Actual 3D mesh creation is in the scope of the project but is not included into the agenda so far.
\end{itemize}
